\section{The affine quadric transfer}\label{sec:vtr}

We use the affine locus $\Sigma\subset\C^{m+1}$, the forms $\lambda,\sigma$ and the exhaustion $\Sigma_\ell$ of
Definition~\ref{def:qdr-curve}, and write $F(z)=(y_j^2-t^2+j)_{j=1}^m$. In \S\ref{sec:vtr-curve}
$m=n-1\ge2$; the transfer of \S\S\ref{sec:vtr-bundle}--\ref{sec:vtr-return} is valid for every $m\ge1$.

\subsection{Sheets and explicit coordinates}\label{sec:vtr-curve}

\begin{lemma}[The affine sheets]\label{lem:vtr-square}
The projection $t:\Sigma\to\C$ has $2^m$ points in every fibre away from $\{\pm\sqrt j:1\le j\le m\}$.
Over $t=\pm\sqrt i$ it has $2^{m-1}$ points, with simple branching in coordinate $y_i$.
Outside a sufficiently large disc its $2^m$ ends are indexed by signs $\epsilon_j\in\{\pm1\}$ and have coordinates
\[
 t=s^{-1},\qquad y_j=\epsilon_j s^{-1}\sqrt{1-js^2},\qquad 0<|s|<1/R_0,
\]
where $R_0=2+\sum_jj$. Each root is the branch equal to one at zero; the inverse is $s=1/t$.
\end{lemma}

\begin{proof}
The independent root signs count the fibres. At a branch value only one root vanishes, and the local equation is
$y_i^2=t^2-i$, with derivative $2t\ne0$ in $t$. The formula at the ends follows by factoring $t^2$ from every
radicand. On $|s|<1/R_0$ all perturbations have modulus less than one, so the roots above and their finite signs
are unambiguous and account for every point.
\end{proof}

\begin{lemma}[Explicit coordinate boxes]\label{lem:vtr-atlas}
\lean{SymplecticPacking.HighDim.ToricPackingStages.ordinaryLocalLift_target,SymplecticPacking.HighDim.ToricPackingStages.ordinaryLocalLift_left,SymplecticPacking.HighDim.ToricPackingStages.ordinaryLocalLift_smooth,SymplecticPacking.HighDim.ToricPackingStages.branchLocalLift_target,SymplecticPacking.HighDim.ToricPackingStages.branchLocalLift_left,SymplecticPacking.HighDim.ToricPackingStages.branchLocalLift_smooth}
\uses{lem:vtr-square}
Let $c=(t_c,y_c)\in\Sigma$.
\begin{enumerate}[label=(\roman*)]
\item If every $y_{c,j}\ne0$, a coordinate near $c$ is $u=t$, with inverse
\[
 t=u,\qquad y_j=y_{c,j}\sqrt{1+(u^2-t_c^2)/y_{c,j}^2}.
\]
\item If $y_{c,i}=0$, a coordinate near $c$ is $v=y_i$, with inverse
\[
 t=t_c\sqrt{1+v^2/t_c^2},\qquad y_i=v,\qquad
 y_j=y_{c,j}\sqrt{1+v^2/y_{c,j}^2}\quad(j\ne i).
\]
Here $t_c^2=i$, and $y_{c,j}^2=i-j\ne0$ for $j\ne i$.
\end{enumerate}
The square roots in (i)--(ii) are the branches equal to one at zero. Restricting a chart to a coordinate square on
which every perturbation of its radicand has modulus less than one gives an injective holomorphic chart.
The inverse coordinates are ambient coordinate projections. Transition functions are compositions of
these projections with the roots above.
\end{lemma}

\begin{proof}
Squaring the coordinates verifies all equations. In (i), at the centre the ratios $y_j/y_{c,j}$ equal one; in (ii)
the same holds for $t/t_c$ and $y_j/y_{c,j}$, $j\ne i$. Requiring positive real part selects these branches uniquely.
The indicated coordinate projection composed with the inverse chart is the identity, while the branch conditions give
the reverse identity. Holomorphicity follows from the binomial series on the prescribed radicand domains.
\end{proof}

\begin{lemma}[A finite rational family on a compact projection domain]\label{lem:vtr-cells}
\uses{lem:vtr-atlas}
Let $R\ge0$ be rational and put
\[
 \eta=\frac1{64},\quad r=\frac1{262144(R+1)},\quad h=\frac r{32},\quad
 N=\lceil R/h\rceil+2.
\]
Consider the grid $q=h(a+ib)$ with $-N\le a,b\le N$. Retain the ordinary centres
\[
 c(q,\epsilon)=\bigl(q,(\epsilon_j\sqrt{q^2-j})_{j=1}^m\bigr)
\]
that pass the rational tests
\[
 |q|^2<(R+\tfrac12)^2,\qquad |q^2-j|^2>(\eta^2/2)^2\quad(1\le j\le m).
\]
Use all signs, ordinary coordinate squares of outer half-width $r$ and inner half-width $r/4$, and the branch charts
centred at every $t_c=\pm\sqrt i$, $y_{c,i}=0$, $y_{c,j}=\pm\sqrt{i-j}$, with outer half-width $1/8$ and inner half-width
$1/32$. The inner open squares cover $\{|t|\le R\}\subset\Sigma$.
Every ordinary closed outer square lies in $\{|t|<R+1\}$. The same holds for the branch squares provided
$i+1/16<(R+1)^2$ for all $i$.
All index tests are comparisons of rational polynomials in $\operatorname{Re}q$, $\operatorname{Im}q$.
\end{lemma}

\begin{proof}
The chosen constants satisfy
\[
 2r\bigl(2(R+1)+2r\bigr)<\eta^2/2,\qquad
 4r(2R+4r)<\eta^2,\qquad 2h\le r/8,\qquad R+h<Nh.
\]
If every $|y_j|\ge\eta$, choose a grid point with $|q-t|<2h\le r/8$. Then
\[
 |q^2-t^2|\le |q-t|(|q|+|t|)<\eta^2/2,
\]
so $|q^2-j|>\eta^2/2$ for each $j$, and the centre passes the tests.
Among its finite root signs there is one for which all ratios to the coordinates of the point have positive real part.
The ordinary chart formula then puts the point in the inner square. The first displayed inequality bounds every
radicand perturbation on the outer square strictly below one; also $|q|+2r<R+1$.
If some $|y_i|<\eta=(1/8)/8$, the separated parameters give $|t|^2$ close to $i$ and $|y_j|^2$ close to $|i-j|$.
The branch formula with the corresponding finite root signs applies. Its perturbations have modulus at most
$4(1/8)^2<1$, and the point is in the inner square. On its closed outer square,
$|t|^2=|i+y_i^2|\le i+1/16<(R+1)^2$.
\end{proof}

\begin{proposition}[Fixed handle data]\label{prop:vtr-handle}
\uses{lem:vtr-atlas, lem:vtr-cells, M-def:hdl-handle-data, def:qdr-model, lem:pre-fs}
On each compact truncation containing the fixed supports, the quadric model has the annular data below. Its primitive operator acts on zero-mass forms with compact support in the affine interior. The annular charts and closed dual form can be
chosen independently of $h$ by the following root formulas.
Put $P(z)=(z+z^{-1})/2$ and $M(z)=(z-z^{-1})/2$. Two annuli, initially defined near $|z|=1$, are
\[
 \begin{split}
 A(z)&=\bigl(\sqrt{1+P(z)^2},\ P(z),\ M(z),\ (i\sqrt{j-1-P(z)^2})_{j=3}^m\bigr),\\
 C(z)&=\bigl(P(z),\ M(z),\ (i\sqrt{j-P(z)^2})_{j=2}^m\bigr).
 \end{split}
\]
Their inverses are respectively $z=y_1+y_2$ and $z=t+y_1$. The transverse cylinder is
$e(s,\theta)=C(\exp(s+2\pi i\theta))$.
Fix $\varepsilon=1/8$ and $w=1/4$, and let
\[
 T_\varepsilon(z)=S\!\left(\frac{|z|^2-(1-\varepsilon)^2}{(1+\varepsilon)^2-(1-\varepsilon)^2}\right),
 \qquad \beta_0=A_*dT_\varepsilon,
\]
with zero extension outside $A$. Then $\beta_0$ is smooth, closed, compactly supported and has period one on the core
of $C$. With $b(\theta)=2S'(2\theta-1/2)$ on $[0,1]$, extended periodically, the fixed smooth circle clock supported in
$[1/4,3/4]$ with integral one, write
\[
 e^*\beta_0=A_0(s,\theta)\,ds+B_0(s,\theta)\,d\theta.
\]
The periodic function
\[
 F(s,\theta)=\int_0^s A_0(v,0)\,dv+\int_0^\theta(B_0(s,v)-b(v))\,dv
\]
satisfies $dF=e^*\beta_0-b(\theta)d\theta$. With
\[
 \chi(s)=S\bigl((7w/8-s)4/w\bigr)S\bigl((7w/8+s)4/w\bigr),\qquad
 \beta=\beta_0-d\bigl(e_*(\chi F)\bigr),
\]
the function $e_*(\chi F)$ is extended by zero. The form $\beta$ equals $b(\theta)d\theta$ on $|s|<5w/8$.
The annular handle has width $w/2=1/8$. The normalization and layer paths use the smaller working half-width
$b_H=w/4=1/16$ and the radial clock $S((s+b_H)/(2b_H))$; $w=1/4$ is the outer cylinder width.

In the product fibre coordinates,
\[
 \Omega_0=\omega_V+d\Gamma_{\mathcal A},\qquad
 \Gamma_{\mathcal A}=(2|p|-1)\vartheta,\qquad
 \vartheta=-d^c\log(1+|t|^2+\textstyle\sum_j|y_j|^2).
\]
The auxiliary area form and the zero-mass primitive operator are obtained from the fixed finite coordinate families
of Lemma~\ref{lem:vtr-cells}, the chart bump partition and local integral primitive of
Construction~\ref{cons:shd-primitive}; their annular correction is the same two-segment formula for $F$.
\end{proposition}

\begin{proof}
The identity $P^2-M^2=1$ verifies the two annulus formulas. On $|z|=1$, $P(z)$ is real in $[-1,1]$ and $M(z)$ is purely
imaginary. Every remaining square-root radicand is positive, including $1+P^2$ and $j-1-P^2$ for $j\ge3$.
Thus the formulas extend to annuli of positive width. The displayed inverses prove injectivity there.
The two core circles meet once: $y_1$ is real on the first and imaginary on the second, hence zero at an intersection;
the positive choice of $t$ then gives $t=1$, $y_2=i$. The radial coordinate crosses the transition in $T_\varepsilon$
with positive orientation at this point. Integrating its derivative gives one, so $\int_{C(|z|=1)}\beta_0=1$.
The derivative of $T_\varepsilon$ vanishes near both boundary circles; extension by zero is smooth and closed.
Closedness gives $\partial_\theta A_0=\partial_s B_0$. Differentiating the two integrals proves the formula for $dF$.
The period-one identity makes $F$ periodic. The cutoff equals one on $|s|\le5w/8$ and vanishes for $|s|\ge7w/8$,
which proves the assertions about $\beta$ and smooth zero extension. All these supports are compact in the affine curve.

The fixed widths are admissible uniformly in the number of quadrics. On $3/4\le|z|\le5/4$ one has
$|P(z)|\le4/3$ and $|\operatorname{Im}P(z)|\le7/24$. Hence $\operatorname{Re}(1+P(z)^2)>0$ and
$\operatorname{Re}(j-1-P(z)^2)>0$ for $j\ge3$, so $A$ is defined on the whole closed annulus $\frac34\le|z|\le\frac54$, which contains the support
$\frac78\le|z|\le\frac98$ of $dT_\varepsilon$.
For $|s|\le1/2$, the inequality $e^{1/2}<7/4$ gives $|P(e^{s+2\pi i\theta})|\le11/8<\sqrt2$;
thus every $j-P^2$, $j\ge2$, has positive real part, and $C$ is defined on that whole closed cylinder.
The chosen widths leave a neighbourhood of the supports inside both coordinate domains.
The finite chart partition and zero-mass local primitives give the remaining scalar handle data by the construction of
Construction~\ref{cons:shd-primitive}. The initial boxes cover the required compact support.
Additional rational subboxes in the connected open affine base connect their overlap graph: an open cover of a connected
space has finite overlap chains between any two nonempty members. Each nonempty overlap contains a rational subbox
with compact closure in both charts. Enumerating these subboxes and their strict inclusion certificates therefore
finds finite chains. The normalized chart bump weights give the partition; the chains give the remaining finite
data of that construction.
Finally $\vartheta=-\lambda$, so $(2|p|-1)\vartheta=(1-2|p|)\lambda$ gives the stated horizontal primitive directly.
\end{proof}

\begin{lemma}[Explicit truncation error]\label{lem:vtr-tail}
\uses{lem:vtr-square, lem:vtr-atlas, lem:qdr-curve}
Put $R_l=2+\sum_{j=1}^m j+l$, $\Sigma_l=\{|t|\le R_l\}\subset\Sigma$, and
$A_l=\int_{\Sigma_l}\sigma$. Then $\Sigma_l$ is a compact connected surface with smooth boundary, and
\[
 0\le2^m-A_l\le\frac{2^m(m+2)}{R_l^2}
 \le\frac{2^m(m+2)}{l+2}.
\]
In particular, for $d>0$ and
\[
 l=\left\lceil\frac{2^{m+6}(m+2)}d\right\rceil,
 \qquad H_l(p)=A_l(1-2|p|),
\]
one has $0\le\Hb(p)-H_l(p)\le d/64$ on $\Delta_h$.
The handle and dual supports of Proposition~\ref{prop:vtr-handle} lie in the interior of $\Sigma_l$.
\end{lemma}

\begin{proof}
The projection $t$ is proper on $\Sigma$, since $|y_j|^2=|t^2-j|\le|t|^2+j$. Its critical points have some $y_j=0$, i.e.\ $t^2=j$,
so the circle $|t|=R_l$ contains no critical value and $\Sigma_l$ has smooth boundary. The interior is connected: analytic
continuation around a small circle about $\sqrt j$ changes the sign of $y_j$ and of no other coordinate, so the monodromy acts
transitively on the $2^m$ sheets over a regular point of the disc. The handle supports lie in $|t|\le\frac53$
(proof of Proposition~\ref{prop:vtr-handle}: $|P|\le\frac43$ on $A$ and $|P|\le\frac{11}8$ on $C$), inside $\Sigma_0$ since $R_l\ge5$.

Let $R\ge R_0=2+\frac{m(m+1)}2$. By Lemma~\ref{lem:vtr-square}, $\{|t|>R\}$ is the disjoint union of $2^m$ punctured discs
$0<|s|<1/R$, on each of which $[1:t:y]=[s:1:v]$ with $F(s)=(s,1,v_1(s),\dots,v_m(s))$, $v_j(s)=\epsilon_j\sqrt{1-js^2}$. Since
$\lambda=d^c\log D$ and $\log D=\log|F(s)|^2-\log|s|^2$ with $\log|s|^2$ harmonic, $\sigma=dd^c\log|F(s)|^2$ there, and its density
with respect to Lebesgue measure in $s$ is $\pi^{-1}(|F|^2|F'|^2-|F^\dagger F'|^2)/|F|^4\le|F'|^2/(\pi|F|^2)$. For $|s|\le1/R$ we have
$|v_j|^2=|1-js^2|\ge1-j/R^2\ge\frac12$ and $|v_j'|^2=j^2|s|^2/|v_j|^2\le2j^2/R^2$, and $\sum_jj^2\le R_0^2\le R^2$, so $|F'|^2\le3$; with
$|F|\ge1$ the density is at most $3/\pi$, and each end contributes at most $3/R^2$. Hence
$\int_{|t|>R}\sigma\le2^m\cdot3/R^2\le2^m(m+2)/R^2$.

On the other hand, by Stokes' theorem $A_{R'}=\int_{|t|=R'}\lambda$ for $\Sigma_{R'}=\{|t|\le R'\}$. On each end,
$\lambda=d^c\log|t|^2+d^c\log|F(s)|^2=\frac{d\arg t}{2\pi}+d^c\log|F(s)|^2$, and $|\partial\log|F|^2|=|F^\dagger F'|/|F|^2=O(|s|)$, so the
second term contributes $O(R'^{-2})$ on the circle of length $2\pi/R'$ in $s$. Thus $A_{R'}\to2^m$, and since $\sigma>0$,
$2^m-A_R=\int_{|t|>R}\sigma\in[0,2^m(m+2)/R^2]$. For $R=R_l\ge l+2\ge2$ we get $R_l^2\ge l+2$, which is the second bound. The ceiling
choice of $l$ gives $2^m-A_l\le d/64$, and $0<1-2|p|\le1$ on $\Delta_h$ gives the profile estimate.
\end{proof}

\subsection{The affine potentials}\label{sec:vtr-bundle}

Write $D=1+|z|^2$ and $u=|W|^2/D^2$. On $\C^{m+1}\times\C^m$ put
\[
 \Phi_b=\log D+b\log(1+u),\qquad \Omega_b=dd^c\Phi_b.
\]

\begin{lemma}[Positivity and the graph forms]\label{lem:vtr-form}
\lean{SymplecticPacking.BranchSmall.VT.bundleForm_pos, SymplecticPacking.BranchSmall.VT.returnForm_pos}
For $0<b<1/2$, $\Omega_b$ is a smooth Kähler form. For $T>0$, its restriction to the graph
$j_T(z)=(z,F(z)/T)$ is
\[
 \omega_T=dd^c\!\left[\log D+b\log\left(1+\frac{|F|^2}{T^2D^2}\right)\right].
\]
The path $\omega_s^{\rm ret}=dd^c[\log D+sb\log(1+|F|^2/D^2)]$, $0\le s\le1$, is also Kähler.
\end{lemma}

\begin{proof}
The identity
\[
 \Phi_b=(1-2b)\log D+b\log(D^2+|W|^2)
\]
exhibits its Hessian as a positive combination of Fubini--Study Hessians.
Indeed $D^2+|W|^2$ is the squared norm of the vector consisting of the ordered quadratic monomials
of $(1,z)$ and the coordinates $W$. Its constant entry is one; its linear $z$ entries and its $W$ entries
separate tangent vectors. The first summand is positive on nonzero base vectors, and the second is positive
on nonzero pure fibre vectors, so the sum is positive definite. The graph identity is pullback by the holomorphic
map $j_T$. The same decomposition with $sb$ proves positivity of the return path.
\end{proof}

\begin{construction}[Radial injection]\label{cons:vtr-injection}
For $0<b<1/2$, let
\[
 \iota_b(z,U)=\left(z,\frac{D(z)\sqrt\pi\,U}{\sqrt{b-|p|}}\right),\qquad
 \iota_b^{-1}(z,W)=\left(z,\sqrt{\frac b\pi}\frac{W}{\sqrt{D(z)^2+|W|^2}}\right).
\]
The first map has domain $|p|<b$ over $\Sigma$, and the second is defined on $\Sigma\times\C^m$.
\end{construction}

\begin{lemma}[The radial pullback]\label{lem:vtr-injection}
\lean{SymplecticPacking.BranchSmall.radial_form}
These maps are mutually inverse smooth maps, and $\iota_b^*\Omega_b=\Omega_0$ on the tube $|p|<b$.
Moreover $u\circ\iota_b=|p|/(b-|p|)$.
\end{lemma}

\begin{proof}
With $|W|^2=D^2\pi|U|^2/(b-|p|)$ we get $u=|p|/(b-|p|)$, $1+u=b/(b-|p|)$ and $|p|=bu/(1+u)$, which verifies the inverse
formulas; both denominators are positive on their domains. On $\{W\ne0\}$, $d^c\log(1+u)=\frac u{1+u}(d^c\log|W|^2-2\lambda)$ and
$d^c\log|W|^2=\operatorname{Im}(W^\dagger dW)/(2\pi|W|^2)$. Along $\iota_b$, $W=c\,U$ with $c>0$ a function of $(z,|p|)$, so
$\operatorname{Im}(W^\dagger dW)=c^2\operatorname{Im}(U^\dagger dU)$ and $b\frac u{1+u}d^c\log|W|^2=|p|\frac{\operatorname{Im}(U^\dagger dU)}{2\pi|U|^2}
=\frac12\operatorname{Im}(U^\dagger dU)=\sum_jp_jd\theta_j$. Hence
\[
 \iota_b^*d^c\Phi_b=\lambda+\sum_jp_jd\theta_j-2|p|\lambda=\sum_jp_jd\theta_j+(1-2|p|)\lambda ,
\]
which is smooth on the axes. Its exterior derivative is $\omega_V+d[(1-2|p|)\lambda]=\Omega_0$.
\end{proof}

\subsection{Complete transport in affine coordinates}\label{sec:vtr-transport}

Let $J_z=dF_z$. At $(z,W)\in X_T=\{F(z)=TW\}$ define $V_T=(\dot z,\dot W)$ by
\begin{equation}\label{eq:vtr-transport}
 J_z\dot z-T\dot W=W,\qquad
 \Omega_b(V_T,k)=0\quad(k\in\ker(J_z,-T)).
\end{equation}
Here $\ell=(J_z,-T)$ is complex linear and onto: at $T=0$ because $J_z$ is onto at points of $\Sigma$ (Lemma~\ref{lem:qdr-curve}),
for $T\ne0$ because of the fibre term. The $\Omega_b$-orthogonal complement of the complex subspace $\ker\ell$ is its orthogonal
complement for the Kähler metric, on which $\ell$ is an isomorphism; so $V_T$ exists, is unique, and depends smoothly on
$(z,W,T)$ on the set $\mathcal X=\{F(z)=TW\}$, which is a smooth submanifold of $\C^{m+1}\times\C^m\times\R$ near $T=0$ and the
graph of $W=F(z)/T$ for $T>0$.

\begin{lemma}[Extension to projective space]\label{lem:vtr-affine-growth}
Let $\overline{\C^{m+1}}=\CP^{m+1}$ with homogeneous coordinates $Z=(Z_*,Z')$, $z=Z'/Z_*$, and hyperplane at infinity
$H_\infty=\{Z_*=0\}$. For $T>0$ and $0\le s\le1$ the forms
\[
 \omega_T=dd^c\Bigl[\log D+b\log\Bigl(1+\frac{|F|^2}{T^2D^2}\Bigr)\Bigr],\qquad
 \omega^{\rm ret}_s=dd^c\Bigl[\log D+sb\log\Bigl(1+\frac{|F|^2}{D^2}\Bigr)\Bigr]
\]
extend to Kähler forms on $\CP^{m+1}$, and the functions $g_T=\frac{|F|^2}{T^2D^2+|F|^2}$ and $\log(1+|F|^2/D^2)$ extend smoothly.
Hence the Moser fields $N^{\rm gr}_T$, $\iota_{N^{\rm gr}_T}\omega_T=\frac{2b}T\,d^cg_T$, and $Y_s$,
$\iota_{Y_s}\omega^{\rm ret}_s=-b\,d^c\log(1+|F|^2/D^2)$, extend to smooth vector fields on $\CP^{m+1}$, smooth in $T>0$ and $s$, and
tangent to $H_\infty$. Their flows over any compact parameter interval are complete on $\CP^{m+1}$ and preserve $\C^{m+1}$.
\end{lemma}

\begin{proof}
Let $F^{\rm hom}_j(Z)=Z_{y_j}^2-Z_t^2+jZ_*^2$ and let $V_2(Z)$ be the vector of ordered quadratic monomials $Z_aZ_b$, so
$|V_2(Z)|^2=|Z|^4$. Put $P_T(Z)=(V_2(Z),F^{\rm hom}(Z)/T)$, a homogeneous quadratic map without nontrivial zero. At $Z=(1,z)$,
$D^2+|F|^2/T^2=|P_T(Z)|^2$, so
\[
 \omega_T=dd^c\bigl[(1-2b)\log|Z|^2+b\log|P_T(Z)|^2\bigr]=(1-2b)\,\omega_{\rm FS}+b\,P_T^*\omega_{\rm FS},
\]
which is a Kähler form on $\CP^{m+1}$ because $1-2b>0$ and the second term is semipositive; likewise
$\omega^{\rm ret}_s=(1-2sb)\omega_{\rm FS}+sb\,P_1^*\omega_{\rm FS}$ with $1-2sb\ge1-2b>0$. The functions
$g_T=|F^{\rm hom}/T|^2/|P_T|^2$ and $\log(1+|F|^2/D^2)=\log(|P_1|^2/|Z|^4)$ are quotients of homogeneous functions of equal
degree with nonvanishing denominators, hence smooth on $\CP^{m+1}$. The Moser fields are therefore smooth on the compact
manifold $\CP^{m+1}$, and their flows are complete.

The involution $\iota(Z_*,Z')=(-Z_*,Z')$ is unitary, $F^{\rm hom}\circ\iota=F^{\rm hom}$ and $|V_2\circ\iota|=|V_2|$, so it preserves all the forms
and functions above, hence the fields (Lemma~\ref{lem:pre-moser}(g)). At a point $x\in H_\infty$, which is fixed by $\iota$, a
$\iota$-invariant vector lies in the $+1$ eigenspace of $d\iota_x$, which is $T_xH_\infty$, since $d\iota_x$ acts by $-1$ on the normal
direction. So the fields are tangent to $H_\infty$, their flows preserve $H_\infty$, and hence its complement $\C^{m+1}$.
\end{proof}

\begin{proposition}[Complete horizontal transport]\label{prop:vtr-transport}
\lean{SymplecticPacking.BranchSmall.VT.transport_trajectory_complete,
SymplecticPacking.BranchSmall.transport_form}
For every point of $X_0=\Sigma\times\C^m$ the field \eqref{eq:vtr-transport} has a unique trajectory to $T=1$.
Its time-one map $\tau$ is smooth and injective and satisfies
$\tau^*(\Omega_b|_{X_1})=\Omega_b|_{X_0}$.
Along each trajectory the normalized fibre norm $u=|W|^2/D^2$ is nonincreasing for $T\ge0$; more precisely,
\[
 \frac{du}{dT}=-\frac{2\pi T}{b}(1+u)^2\Omega_b(V_T,iV_T)\le0.
\]
\end{proposition}

\begin{proof}
\emph{Constraint.} Along a trajectory, $\frac d{dT}(F(z)-TW)=J_z\dot z-W-T\dot W=0$, so the trajectory stays in $X_T$.

\emph{The graph equation.} For $T>0$ identify $X_T$ with $\C^{m+1}$ by $j_T(z)=(z,F(z)/T)$; then $j_T^*\Omega_b=\omega_T$
(Lemma~\ref{lem:vtr-form}). Write $V_T=dj_T(\dot z)-(0,W/T)$, using $\dot W=(J_z\dot z-W)/T$. For $v\in\C^{m+1}$, $dj_T(v)\in\ker\ell$, so
$\omega_T(\dot z,v)=\Omega_b((0,W/T),dj_T(v))$. The fibre circle action is generated by $X=2\pi i(0,W)$ and preserves $\Phi_b$, and
$d^c\Phi_b(X)=bu/(1+u)$; hence $\iota_X\Omega_b=-d(bu/(1+u))$. Using $J$-invariance of $\Omega_b$ and holomorphy of $j_T$,
$\Omega_b((0,W),dj_T(v))=\Omega_b(i(0,W),dj_T(iv))=-\frac1{2\pi}d\bigl(b\tfrac u{1+u}\bigr)(dj_T(iv))$, and for real $\phi$,
$d\phi(iv)=-4\pi\,d^c\phi(v)$. Since $u\circ j_T=|F|^2/(T^2D^2)$, so that $\frac u{1+u}\circ j_T=g_T$,
\[
 \iota_{\dot z}\omega_T=\frac{2b}T\,d^cg_T ,
\]
that is, $\dot z=N^{\rm gr}_T(z)$ is the field of Lemma~\ref{lem:vtr-affine-growth}. (Equivalently $\dot z=-r_TN_T/T$ with
$\iota_{N_T}\omega_T=-b\,d^c\log(1+u\circ j_T)$ and $r_T=2/(1+u\circ j_T)$.)

\emph{Completeness.} Let $x\in X_0$. Since $V_T$ is smooth on $\mathcal X$, the trajectory exists on $[0,\delta]$ for some
$\delta\in(0,1)$. For $T\in[\delta,1]$ it is, in the graph coordinate, an integral curve of $N^{\rm gr}_T$, whose flow is complete in
$\C^{m+1}$ by Lemma~\ref{lem:vtr-affine-growth}; so it extends to $[0,1]$, and it is unique. The time-$T$ maps are restrictions of
the flow of the smooth field $V_T+\partial_T$ on $\mathcal X$, hence smooth and injective (backward uniqueness).

\emph{Norm.} By the circle identity above, $\frac d{dT}\frac{bu}{2\pi(1+u)}=\Omega_b(V_T,i(0,W))$. Since $\ell(i(0,W))=-iTW$ and
$\ell(iV_T)=iW$, the vector $i(0,W)+T\,iV_T$ lies in $\ker\ell$, so $\Omega_b(V_T,i(0,W))=-T\,\Omega_b(V_T,iV_T)$. This gives the
displayed derivative, which is $\le0$ for $T\ge0$ because $\Omega_b(V,iV)$ is the Kähler norm of $V$.

\emph{Pullback.} Let $\Phi(\sigma,T)$ be the transported image of a smooth family of initial points. Its tangent vectors are
$(\kappa,t)$ with $\ell(\kappa)=tW$, and $\partial_T\Phi=(V_T,1)$. Writing $\kappa=tV_T+(\kappa-tV_T)$ with $\kappa-tV_T\in\ker\ell$
gives $\Omega_b(V_T,\kappa)=0$, so $\iota_{\partial_T}\Phi^*\Omega_b=0$. As $\Phi^*\Omega_b$ is closed, Cartan's formula gives
$\partial_T(\Phi_T^*\Omega_b)=0$, which is the pullback identity.
\end{proof}

\subsection{Moser return and the ball chart}\label{sec:vtr-return}

\begin{proposition}[Affine return]\label{prop:vtr-return}
\lean{SymplecticPacking.BranchSmall.VT.return_control_all, SymplecticPacking.BranchSmall.returnInv_form}
Let
\[
 \omega_s^{\rm ret}=dd^c\!\left[\log D+sb\log\left(1+\frac{|F|^2}{D^2}\right)\right],\qquad
 \iota_{Y_s}\omega_s^{\rm ret}=-b\,d^c\log\left(1+\frac{|F|^2}{D^2}\right).
\]
The smooth field $Y_s$ has a flow through every affine point for all $s\in[0,1]$, in both directions.
If $\psi_s$ starts at $s=0$, then $\psi_s^*\omega_s^{\rm ret}=dd^c\log D$.
\end{proposition}

\begin{proof}
By Lemma~\ref{lem:vtr-affine-growth}, $Y_s$ is a smooth field on $\CP^{m+1}$ tangent to $H_\infty$, so its flow is defined for all
$s$ in both directions and preserves $\C^{m+1}$. Since
$\partial_s\omega^{\rm ret}_s=d\bigl(b\,d^c\log(1+|F|^2/D^2)\bigr)$, Cartan's formula gives
$\frac d{ds}\psi_s^*\omega^{\rm ret}_s=\psi_s^*d\bigl(\iota_{Y_s}\omega^{\rm ret}_s+b\,d^c\log(1+|F|^2/D^2)\bigr)=0$, so
$\psi_s^*\omega_s^{\rm ret}=\omega^{\rm ret}_0=dd^c\log D$. Finally $j_1^*\Omega_b=\omega^{\rm ret}_1$ by Lemma~\ref{lem:vtr-form}.
\end{proof}

\begin{theorem}[The affine quadric transfer]\label{thm:vtr-transfer}\label{thm:qdr-transfer}
\lean{SymplecticPacking.BranchSmall.vectorTransfer_spec}
\uses{lem:vtr-injection, prop:vtr-transport, prop:vtr-return}
For $0<b<1/2$, the map
\[
 \mathcal T_b(z,U)=\mathrm{ch}\bigl(\psi_1^{-1}(\mathrm{pr}(\tau(\iota_b(z,U))))\bigr),\qquad
 \mathrm{ch}(z)=\frac{z}{\sqrt{\pi(1+|z|^2)}},
\]
is defined on the whole tube $U_b=\{(z,U)\in\Sigma\times\C^m:|p(U)|<b\}$ and is a smooth symplectic embedding of
$(U_b,\Omega_0)$ into $(\Bo(1),\omega_0)$, for every $m\ge1$. It agrees with $\mathrm{ch}$ on the zero section. For every compact
$\mathcal K\subset U_b$ there is $m_{\mathcal K}>0$ with $\mathcal T_b(\mathcal K)\subset B(1-m_{\mathcal K})$. For $0<h<b$ the map is
defined on a neighbourhood of every finite model $V_{\Delta_h}(E_\ell)$; it depends on $b$ and the fixed curve, and not on
$h$, $\ell$ or $\mathcal K$.
\end{theorem}

\begin{proof}
By Lemma~\ref{lem:vtr-injection}, Proposition~\ref{prop:vtr-transport} and Proposition~\ref{prop:vtr-return},
\begin{align*}
 \mathcal T_b^*\omega_0&=\iota_b^*\tau^*\mathrm{pr}^*(\psi_1^{-1})^*\mathrm{ch}^*\omega_0
 =\iota_b^*\tau^*\mathrm{pr}^*\omega^{\rm ret}_1\\
 &=\iota_b^*\tau^*(\Omega_b|_{X_1})=\iota_b^*\Omega_b=\Omega_0 ,
\end{align*}
using $\mathrm{ch}^*\omega_0=dd^c\log D$ (Lemma~\ref{lem:pre-fs}),
$(\psi_1^{-1})^*dd^c\log D=\omega^{\rm ret}_1$, and $\mathrm{pr}^*\omega^{\rm ret}_1=\mathrm{pr}^*j_1^*\Omega_b=\Omega_b|_{X_1}$ since
$j_1\circ\mathrm{pr}=\mathrm{id}$ on $X_1$. Each factor is injective: $\iota_b$ has an inverse, $\tau$ and $\psi_1^{-1}$ are time maps of flows,
$\mathrm{pr}$ is injective on the graph $X_1$, and $\mathrm{ch}$ is a bijection onto $\Bo(1)$ (Lemma~\ref{lem:pre-fs}(b)). A
symplectic map between manifolds of equal dimension $2n$ is an immersion, so $\mathcal T_b$ is an injective local
diffeomorphism, hence an open embedding, with values in $\Bo(1)$ because $\pi|\mathrm{ch}(z)|^2=1-D^{-1}<1$. At $U=0$, $W=0$ and
$V_T=0$, so $\tau(z,0)=(z,0)$; and $Y_s$ vanishes on $\Sigma=\{F=0\}$, where $d|F|^2=0$; so $\mathcal T_b(z,0)=\mathrm{ch}(z)$. A compact
source has compact image in the open ball, at positive distance from its boundary sphere. For $0<h<b$ the closed set
$V_{\Delta_h}(E_\ell)$ lies in the open tube $U_b$, and no choice in the construction involves $h$, $\ell$ or $\mathcal K$.
\end{proof}
